% CoRL 2026 PhysWM workshop paper, 4 pages plus references.
% No CoRL template is checked into this repo. When the workshop's corl_2026.sty (or the
% PhysWM variant) is available, replace the preamble block below with
%   \documentclass{article}\usepackage{corl_2026}  (plus \usepackage[final]{corl_2026} for camera-ready)
% and keep everything from \title onward. The body uses only booktabs, graphicx, amsmath,
% xcolor, multirow, natbib and hyperref, all of which the CoRL template loads or tolerates.
% The CoRL text block (5.5in) is narrower than this preamble's (6.7in); expect the draft to run
% about 15% longer there and cut accordingly.
%
% Title alternatives (pick one, delete the others):
%   A. What to Start From: Pretraining Exposure, Not Architecture, Predicts Small-Budget Doom World Models
%   B. Exposure Beats Architecture: Choosing a Warm Start for Small-Budget Doom World Models
%   C. Start From What Saw Your Latents: A Controlled Warm-Start Study on an Open Doom Benchmark
\documentclass[10pt]{article}
\usepackage[margin=0.9in]{geometry}
\usepackage{booktabs,graphicx,amsmath,microtype,xcolor,multirow}
\usepackage[numbers,sort&compress]{natbib}
\usepackage[colorlinks,linkcolor=blue,citecolor=blue,urlcolor=blue]{hyperref}
\newcommand{\todo}[1]{\textcolor{red}{[TODO: #1]}}
\newcommand{\tbd}{\textcolor{red}{[tbd]}}  % compact placeholder for pending table cells
\newcommand{\figplaceholder}[2]{\IfFileExists{#1}{\includegraphics[width=#2]{#1}}{\fbox{\parbox[c][2.6cm][c]{\dimexpr#2-2\fboxsep-2\fboxrule\relax}{\centering\todo{figure \texttt{\detokenize{#1}}}}}}}

\title{What to Start From: Pretraining Exposure, Not Architecture,\\ Predicts Small-Budget Doom World Models}
\author{Rohan Nagabhirava \and Keerthana Chirumamilla \and Changliu Liu\\ Carnegie Mellon University}
\date{}
\begin{document}
\maketitle

\begin{abstract}
A lab that wants a world model of one environment usually has a small dataset, a few GPUs, and a shelf of pretrained diffusion checkpoints to start from. We release the first open, reproducible Doom world-model benchmark: 850 lossless Freedoom episodes on 17 maps (1.05M verified decision frames), a seeded evaluation corpus with held-out maps, persistence and autoencoder references, an inverse-dynamics judge, drift by horizon, FVD, and episode-level confidence intervals. On it we run a controlled warm-start study under one recipe (90k updates, batch 32) across five pretrained starts: the ImageNet DiT-XL/2, PixArt-$\alpha$, the Stable Diffusion 1.4 U-Net, UniDiffuser, and a video-pretrained transformer. The standard ImageNet DiT is the weakest start on every metric; starts that saw broad data in the target latent space do better whether they are U-Nets or transformers. We also find that by 64 rollout steps every model is perceptually worse than a frozen seed frame, so long-horizon fidelity must be read against persistence.
\end{abstract}

\section{Introduction}
GameNGen~\citep{valevski2024gamengen} and MultiGen~\citep{po2026multigen} simulate Doom with a Stable Diffusion U-Net trained on hundreds of millions of frames; the wider world-model literature has moved to diffusion transformers fine-tuned from video foundation models~\citep{he2025matrixgame2,oasis2024}. A lab with one environment, one million frames, and a few 48\,GB GPUs can follow neither recipe at scale and must first answer a narrower question: \emph{at a fixed small budget, which pretrained checkpoint should the world model start from?} Published comparisons cross papers, and so cross data, resolution, context, objective, and compute at once. We use Doom as a testbed for this general question, not as the goal: it has a public reference recipe, a scripted agent that plays many maps, and a rendering path we control down to the engine tic, so everything but the starting checkpoint can be held fixed. The robot lab fine-tuning a foundation model on a small in-domain corpus is the reader we have in mind.

\textbf{Contribution 1, the vehicle: an open Doom benchmark.} GameNGen and MultiGen released neither data nor weights; PlayGen~\citep{playgen2024} released code but no comparable corpus. We release a redistributable Freedoom corpus (850 episodes, 17 maps, every engine tic lossless, 1.05M verified decision frames), a seeded evaluation corpus with held-out maps, and a protocol with persistence and autoencoder references, a blur control, episode-level bootstrap intervals, an inverse-dynamics judge validated against its two-frame shortcut, drift by horizon, and FVD (Section~\ref{sec:bench}).

\textbf{Contribution 2, the headline: a controlled warm-start study.} Under one recipe (Section~\ref{sec:setup}) we adapt five starts spanning two architectures and three orders of magnitude of pretraining exposure: DiT-XL/2 from ImageNet (1.3M images), PixArt-$\alpha$ (25M), the SD~1.4 U-Net (LAION, about 2B), UniDiffuser (LAION-2B, a transformer), and, exploratory and in its own autoencoder, the video-pretrained SkyReels-V2 transformer. The ImageNet DiT is the weakest start we measured; PixArt-$\alpha$, a transformer with 20 times its exposure, reaches a lower held-out loss than the U-Net. We state this as consistent with exposure in the target latent space predicting quality, not as causal: the rows also differ in size, conditioning path, and objective.

\textbf{Contribution 3, supporting: map transfer.} The U-Net wins perceptual quality on every training map, but on the two held-out maps the transformers lead PSNR, match LPIPS, and degrade less (1.54 versus 2.23\,dB); a corpus on 15 further maps is being scored to test this beyond two maps. A fourth finding is about evaluation: at horizon 64 the frozen seed frame beats every model on LPIPS, so a long-horizon PSNR lead is not by itself evidence of better dynamics.

\section{Benchmark}
\label{sec:bench}
\textbf{Data.} We run ViZDoom~\citep{kempka2016vizdoom} deathmatch with Freedoom assets~\citep{freedoom} on the 17 maps of the Arnold WAD, 8 bots, 150 game-seconds per episode, 50 episodes per map, at 320$\times$240 with HUD, weapon, and crosshair. The player is the unmodified 2017 Arnold agent~\citep{lample2017arnold}, which MultiGen also used; it decides every 4 tics from 29 discrete actions. Every engine tic (35 per second) is stored losslessly with the action applied from that tic, the button vector, HUD state, and player pose: 850 episodes, 4.21M tics, 33.4 hours, 202\,GB. Freedoom is BSD-licensed, so the corpus is redistributable; the public alternatives are JPEG-compressed, 60$\times$80, or single-arena.

\textbf{Verified transitions.} The prediction unit is one frame per agent decision (stride 4): per tic, copy-last already scores 26.2\,dB, against 19.4\,dB per decision. Deaths shift the decision phase, so 62\% of decisions start off the tic-modulo-4 grid and a naive stride-4 slice straddles two actions in 39\% of windows. We cut windows at recorded decision boundaries and keep transitions whose action was held for exactly 4 tics: 1,054,647 verified decision frames (99.4\%) in 7,296 chains of median length 108, so every target carries one verified action.

\textbf{Splits and evaluation corpus.} Maps 16 and 17 are held out entirely; 10\% of episodes on maps 1--15 monitor training (675 / 75 / 100 episodes). Every reported number comes from a separate seeded corpus (the recorder seeds every RNG stream from corpus and episode id; re-recording reproduces every tic and frame hash) of 60 seen-map episodes (4 per map 1--15) and 20 unseen-map episodes (10 each on maps 16, 17) that touched no fitted component, decoder and judge included. A second unseen corpus on 15 further maps is being recorded \todo{name the map set and episode count}.

\textbf{Metrics and references.} Teacher-forced: 32 real context frames and the action in, one sampled frame out, scored against the raw frame on 2,048 windows per corpus with PSNR and LPIPS (AlexNet)~\citep{zhang2018lpips}. Two references bracket each row: \emph{copy-last}, the last context frame, which a model must beat to have learned dynamics, and the \emph{VAE ceiling}, the true latent decoded, which no latent model can exceed. Autoregressive: from 32 real seed frames and the recorded actions, 64 decision frames fed back, 256 rollouts on the seen corpus, identical trajectories and actions for every row (checksum-verified). We report PSNR and LPIPS per horizon, FVD~\citep{unterthiner2018fvd} on 16- and 32-frame clips, and \emph{copy-seed}, the last seed frame repeated, as the persistence reference. A blur control rescores predictions after Gaussian blur ($\sigma$ up to 4\,px) to test whether a PSNR lead can be manufactured by smoothing. Uncertainty is an episode-level bootstrap (95\%), since windows within an episode are correlated.

\textbf{Judge.} Action following is read back by an inverse-dynamics model (IDM) in the structure of VPT's~\citep{baker2022vpt}, as in MineWorld: a convolutional encoder over a bidirectional window of $K{=}8$ latents and a 4-layer transformer, 11.4M parameters, trained on real transitions from the training split. On held-out real windows it reaches 82.9\% top-1 over 29 actions (majority 23.2\%); a $K{=}2$ version reaches 70.1\%, so the window adds 13 points rather than inventing the signal. On the real counterparts of the evaluation rollouts it scores 0.847, the reference for every generated-rollout number.

\section{Setup}
\label{sec:setup}
\textbf{Shared recipe.} Frames are padded to 320$\times$256 and encoded once with the frozen SD KL-f8 encoder (\texttt{sd-vae-ft-mse}, scale 0.18215) to $4\times32\times40$ latents, the space four of the five starts were pretrained in. The decoder is fine-tuned on training frames with MSE $+\,0.1\times$LPIPS, as in GameNGen, because the frozen decoder reconstructs Doom at only 23.7\,dB (HUD rows 17.9\,dB); fine-tuned it reaches 28.3\,dB / 0.051 and decodes every number here. The $L{=}32$ context latents are stacked on channels with the noisy target (132 channels); the conditioning action is the verified action of the transition into the target; the context is corrupted in training with Gaussian noise at a level drawn uniformly up to 0.7, discretized into 10 buckets fed as a conditioning id~\citep{valevski2024gamengen}, and clean at inference. Every start predicts the velocity $v$~\citep{salimans2022progressive} under a 1,000-step linear schedule, learned-variance heads removed. AdamW, learning rate $5\times10^{-5}$, 2,000 warmup steps, no weight decay, clipping at 1.0, global batch 32, 90k updates (2.88M target presentations), bf16 autocast over fp32 master weights, seed 0. Sampling is DDIM~\citep{song2021ddim}, 50 steps, $\eta{=}0$. We report the 90k live weights; the EMA grid is in Appendix~\ref{app:ema}. \todo{confirm the checkpoint policy (final 90k versus best held-out) against the evaluation logs.}

\textbf{Warm starts} (Table~\ref{tab:starts}). Each start keeps its native conditioning pathway, since forcing adaLN onto a U-Net or cross-attention onto a DiT would test a modification, not the checkpoint. The input projection is inflated with the pretrained kernel on the four target channels and zeros on the 128 context channels, so step 0 equals the pretrained model applied to the noisy target; positional tables are regenerated for the $16\times20$ grid. DiT-XL/2~\citep{peebles2023dit} takes action and noise bucket as adaLN-Zero tables; PixArt-$\alpha$~\citep{chen2023pixart} as two learned tokens through its caption projection with timestep on adaLN-single; the SD~1.4 U-Net~\citep{rombach2022ldm} as one cross-attention token and a class embedding; UniDiffuser~\citep{bao2023unidiffuser} \todo{conditioning design}; SkyReels-V2~\citep{chen2025skyreels} \todo{action path}. The video row is exploratory: Wan-VAE latents (16 channels, $4\times8\times8$), $L{=}8$, flow matching, 10k updates of batch 8, so 80k presentations, 36 times fewer than the other rows.

\begin{table}[t]\centering\footnotesize
\caption{Warm starts under the shared recipe. Params are the denoiser alone; exposure is the disclosed pretraining image count (SD~1.4: about 1.17M steps at batch 2048 over LAION subsets). Updates/s on one A6000 at global batch 32 is a proxy for adaptation compute \todo{replace with measured card-hours}.}\label{tab:starts}
\begin{tabular}{llllc}\toprule
Start & Type & Params & Pretraining exposure & Updates/s \\\midrule
DiT-XL/2 ImageNet & DiT, adaLN-Zero & 675M & ImageNet-1k, 1.3M images & 1.26 \\
PixArt-$\alpha$ 512 & DiT, adaLN-single + x-attn & 612M & 25M images (ImageNet, SAM, aesthetic) & 0.68$^\dagger$ \\
SD 1.4 U-Net & U-Net, x-attn & 860M & LAION subsets, $\sim$2B pairs & 0.85 \\
UniDiffuser & U-ViT, token sequence & 952M & LAION-2B, -HR, -aesthetics & \tbd \\
SkyReels-V2 DF 1.3B & video transformer, own VAE & 1.44B & video, count not disclosed & 0.17$^\ddagger$ \\\bottomrule
\end{tabular}\\[2pt]
{\footnotesize $^\dagger$measured while sharing the card with another run. $^\ddagger$batch 8 with gradient checkpointing.}
\end{table}

\section{Results}
\label{sec:results}
\begin{table}[t]\centering\footnotesize
\caption{Main results, seed 0, live weights, fresh evaluation corpus. Teacher-forced PSNR (dB) / LPIPS against lossless frames, 2,048 windows per corpus; rollouts 256$\times$64 on seen maps. Copy-seed is the persistence reference for the rollout columns. Second seed, EMA, and per-horizon tables are in the appendix.}\label{tab:main}
\begin{tabular}{lcccccccc}\toprule
 & \multicolumn{2}{c}{Seen maps} & \multicolumn{2}{c}{Unseen maps} & \multicolumn{2}{c}{Rollout @64} & & \\
 & PSNR$\uparrow$ & LPIPS$\downarrow$ & PSNR$\uparrow$ & LPIPS$\downarrow$ & PSNR$\uparrow$ & LPIPS$\downarrow$ & FVD$_{16/32}\downarrow$ & IDM$\uparrow$ \\\midrule
copy-last / copy-seed & 19.41 & -- & 18.50 & -- & 17.86 & 0.510 & -- & -- \\
VAE ceiling / real clips & 28.61 & 0.051 & 26.36 & 0.070 & -- & -- & -- & 0.847 \\\midrule
DiT-XL/2 ImageNet & 21.06 & 0.311 & 19.52 & 0.450 & 18.30 & 0.553 & 232 / 481 & 0.492 \\
PixArt-$\alpha$ & \tbd & \tbd & \tbd & \tbd & \tbd & \tbd & \tbd & \tbd \\
SD 1.4 U-Net & \textbf{21.36} & \textbf{0.270} & 19.14 & 0.446 & 16.03 & 0.566 & \textbf{184 / 356} & \textbf{0.509} \\
UniDiffuser & \tbd & \tbd & \tbd & \tbd & \tbd & \tbd & \tbd & \tbd \\
SkyReels-V2 DF (own VAE) & \tbd & \tbd & \tbd & \tbd & \tbd & \tbd & \tbd & \tbd$^\S$ \\\bottomrule
\end{tabular}\\[2pt]
{\footnotesize $^\S$the video row's rollouts pass through the SD encoder to reach the judge; read against its own real reference.}
\end{table}

\begin{figure}[t]\centering
\figplaceholder{figures/drift_psnr.pdf}{0.48\linewidth}\hfill
\figplaceholder{figures/drift_lpips.pdf}{0.48\linewidth}
\caption{Rollout PSNR (left) and LPIPS (right) by horizon on the seen corpus, 256 rollouts, with the copy-seed persistence reference and episode-bootstrap bands. Every model crosses below copy-seed on LPIPS before horizon 64. \todo{generate from rollout\_eval outputs; add PixArt and UniDiffuser.}}\label{fig:drift}
\end{figure}

\textbf{Exposure, not architecture.} Held-out $v$-loss at 90k orders the finished rows PixArt-$\alpha$ 0.2030 $<$ U-Net 0.2043 $<$ DiT-XL/2 0.2139 (the second DiT seed also ends at 0.2139). The ImageNet DiT is the weakest start on every column of Table~\ref{tab:main} where a comparison exists: the U-Net leads it by 0.041 LPIPS on seen maps (0.270 [0.261, 0.280] versus 0.311 [0.301, 0.320], intervals disjoint) and by 48 FVD$_{16}$, while seen PSNR overlaps (21.36 [21.01, 21.73] versus 21.06 [20.74, 21.38]). PixArt-$\alpha$ shares the DiT's 28-block, 1152-wide skeleton and latent space but saw 20 times more images, and its loss is below the U-Net's despite an 80-fold smaller corpus \todo{fill the PixArt row and rewrite this sentence against its pixel and rollout metrics}. This is consistent with the useful quantity in a warm start being how much of the target latent space it has seen, with architecture second order at this budget. It is not a controlled test of that reading: rows differ in size (612M to 860M), conditioning path, and pretraining objective, and equal updates are not equal compute. UniDiffuser, a LAION-scale transformer in the same latent space, separates architecture from exposure most directly \todo{UniDiffuser result: does it land with the U-Net or with PixArt?}. The video row is a different autoencoder with 36 times fewer presentations; we report it for completeness, not as a sixth point on the exposure axis \todo{SkyReels row, one sentence}.

\textbf{Map transfer.} Table~\ref{tab:transfer} breaks the teacher-forced numbers down by map. The U-Net wins LPIPS on all 15 training maps and PSNR on 13. On the two held-out maps the order flips: both DiT seeds lead PSNR (19.52 [19.32, 19.73] and 19.59 [19.40, 19.80] versus 19.14 [18.90, 19.39], intervals disjoint) and LPIPS is a wash (0.450 / 0.442 versus 0.446); the seen-to-unseen PSNR drop is 1.54\,dB for the DiT and 2.23\,dB for the U-Net. Two maps cannot carry a generalization claim; we state this as consistent with the transformer starts transferring better and defer the claim to the 15-map corpus \todo{fill the 15-map column; upgrade or retract the wording accordingly}.

\begin{table}[t]\centering\footnotesize
\caption{Map transfer, teacher-forced, live weights. Wins count training maps on which the U-Net beats DiT seed 0. Degradation is seen minus unseen PSNR.}\label{tab:transfer}
\begin{tabular}{lcccccc}\toprule
 & \multicolumn{2}{c}{U-Net wins (15 maps)} & \multicolumn{2}{c}{Unseen 16--17} & Degr. & 15 further maps \\
 & PSNR & LPIPS & PSNR & LPIPS & (dB) & PSNR / LPIPS \\\midrule
DiT-XL/2 s0 & \multirow{3}{*}{13 / 15} & \multirow{3}{*}{15 / 15} & 19.52 & 0.450 & 1.54 & \tbd \\
DiT-XL/2 s1 & & & 19.59 & 0.442 & 1.62 & \tbd \\
SD 1.4 U-Net & & & 19.14 & 0.446 & 2.23 & \tbd \\
PixArt-$\alpha$ & \tbd & \tbd & \tbd & \tbd & \tbd & \tbd \\\bottomrule
\end{tabular}
\end{table}

\textbf{Persistence.} Figure~\ref{fig:drift} carries a result about evaluation, not about a model. At horizon 64 the frozen seed frame scores 17.86\,dB / 0.510 LPIPS; the finished models score 0.550 to 0.566 LPIPS, all worse than persistence, and only DiT seed 0 exceeds copy-seed on PSNR (by 0.44\,dB; seed 1 is 0.18\,dB below, the U-Net 1.83\,dB below). The DiT's horizon-64 PSNR lead over the U-Net is real ($-2.28$\,dB [$-2.88$, $-1.70$] for U-Net minus DiT s0) and repeats on the second seed (1.65\,dB), and it is not smoothing: blur never raises PSNR for any model (Appendix~\ref{app:blur}). The horizon-32-to-64 drop is heavy-tailed for every model (25\% to 43\% of rollouts lose more than 2\,dB), not a few bad clips. \todo{motion audit pending; keep exactly one of the next two sentences.} \todo{A: a motion-magnitude audit finds the DiT's frame-to-frame change comparable to the real trajectories', so the lead is not reduced motion either; it remains a pixel-fidelity result and does not establish more accurate dynamics.} \todo{B: a motion-magnitude audit finds the DiT's frame-to-frame change below the real trajectories', so the lead is at least partly a model that moves less, which persistence rewards; it does not establish more accurate dynamics.} Either way, a long-horizon PSNR lead reported without a persistence reference will be over-read.

\subsection*{What this suggests for small-budget world models}
Doom stands in for any single environment with a small in-domain corpus. As conditional guidance rather than law, our measurements suggest:
\begin{itemize}\setlength{\itemsep}{1pt}
\item \emph{Choose the warm start by how much of your latent space it has seen, not by its architecture.} The default ImageNet DiT was the weakest start on every metric; the two starts with broad exposure in the same latent space beat it, U-Net and transformer alike.
\item \emph{Architecture is second order at this budget.} If the strongest checkpoint in your latent space is a U-Net, a U-Net is a reasonable default; a transformer earns its place only with comparable exposure behind it.
\item \emph{Report long-horizon rollouts against a persistence reference.} Every model here fell below the frozen seed frame on LPIPS by 64 steps; without copy-seed in the plot, a PSNR lead reads as dynamics when it may be fidelity to a scene that barely changed.
\item \emph{Measure transfer to unseen scenes.} Seen-scene metrics overstated the U-Net's advantage in our data: it won every training map on LPIPS and lost PSNR on both held-out maps \todo{confirm or revise once the 15-map corpus is scored}.
\end{itemize}

\section{Related work}
GameNGen~\citep{valevski2024gamengen} fine-tunes the SD~1.4 U-Net with a 64-frame channel-stacked context and noise-augmented history for 700k updates at batch 128 on 128 TPU-v5e over 900M frames (29.43\,dB / 0.249 LPIPS, FVD 114 / 186), releasing neither data nor weights. MultiGen~\citep{po2026multigen} keeps the U-Net and adds an explicit map-and-pose memory for level-conditioned and multiplayer Doom on 100 generated maps, disclosing no hyperparameters or release. PlayGen~\citep{playgen2024} is the first DiT for Doom: a 131M transformer in its own VAE with a diffusion-forcing recurrent state, trained on 200M transitions at 128$\times$128 and reported by rollout length with the action-accuracy judge our IDM follows. DIAMOND~\citep{alonso2024diamond} trains pixel-space EDM U-Nets from scratch on Atari and CS:GO on one RTX 4090 and releases playable models. Oasis~\citep{oasis2024} releases a 500M from-scratch DiT with diffusion forcing for Minecraft, with inference code but no training details or evaluation. MineWorld~\citep{mineworld2025} is a 300M to 1.2B autoregressive transformer over VQ tokens on the VPT corpus and validates an IDM controllability metric against human ratings. Matrix-Game 2.0~\citep{he2025matrixgame2} fine-tunes SkyReels-V2-I2V-1.3B on 1,200 hours of action-labeled video for 120k updates at batch 256 and distills it to 3-step streaming, the pattern our fifth row borrows. Nano World Models~\citep{huang2026nanowm} is a controlled codebase at our scale (40M to 830M transformers in SD KL-f8 latents) that finds $\epsilon$-prediction worst and prices action-injection choices; it varies scale and injection, not the warm start, the axis this paper adds.

\section{Limitations and conclusion}
One recipe at one budget: a per-start learning rate or step count could reorder rows, and equal updates are not equal compute. The exposure reading is correlational across five starts that also differ in size, conditioning path, and objective. Map transfer rests on two held-out maps until the 15-map corpus lands; the video row is exploratory. Sampling is 50-step DDIM at about 2 frames per second on one A6000, so nothing here is real time, and the models hold no memory beyond 32 frames. The agent is one scripted policy, not human play, and the judge reads actions from frames and cannot certify causal control. What the benchmark and study do establish is useful to a lab with one environment and a few GPUs: start from the checkpoint that has seen the most of your latent space, expect architecture to matter less than that at this scale, and never read a long-horizon curve without persistence drawn on it.

\bibliographystyle{plainnat}
\bibliography{refs}

\appendix
\section{Second seed}\label{app:seed}
DiT-XL/2 ImageNet seed 1 under the identical recipe: seen 21.21 / 0.307, unseen 19.59 / 0.442, rollout PSNR@8/16/32/64 19.86 / 19.09 / 18.65 / 17.68, LPIPS@8/64 0.436 / 0.550, IDM 0.476, FVD 198 / 407, held-out $v$-loss 0.2139 (unrounded 0.213875 versus seed 0's 0.213941; distinct checkpoints by md5). Seeds agree to 0.15\,dB and 0.004 LPIPS on seen-map teacher forcing and differ by 0.62\,dB at horizon 64 and 74 FVD$_{32}$: seed variance is small for the teacher-forced ordering and material for long-horizon numbers.

\section{Live versus EMA}\label{app:ema}
Seen-map teacher-forced, live (EMA 0.9999, fp32): DiT s0 21.06 / 0.311 (21.06 / 0.306); DiT s1 21.21 / 0.307 (21.10 / 0.304); U-Net 21.36 / 0.270 (21.67 / 0.250). EMA helps the U-Net by 0.31\,dB and 0.020 LPIPS and barely moves the DiTs; the main table is all-live so no row is selected by whichever weights favor it. \todo{PixArt, UniDiffuser rows; unseen and rollout EMA columns.}

\section{Per-horizon tables}\label{app:horizon}
Rollout PSNR@8/16/32/64 and LPIPS@8/64, seen corpus, live weights: DiT s0 19.66 / 18.93 / 18.51 / 18.30, 0.453 / 0.553; DiT s1 19.86 / 19.09 / 18.65 / 17.68, 0.436 / 0.550; U-Net 19.55 / 18.66 / 17.73 / 16.03, 0.421 / 0.566. Median horizon-32-to-64 drop $-0.03$ / $-0.54$ / $-1.20$\,dB; 25 / 36 / 43\% of rollouts drop more than 2\,dB; trimming the worst 10\% shifts every mean by 0.6 to 0.7\,dB. \todo{full table with episode-bootstrap intervals at every horizon; PixArt, UniDiffuser, SkyReels rows.}

\section{Historical grid-label checkpoints}\label{app:grid}
Before the verified-transition encoding, both backbones were trained on windows cut on the tic-modulo-4 grid, where 39\% of windows straddled two actions. Scored on the same fresh corpus with the same suite (unmatched steps, historical reference only): DiT at 75k, seen 20.18 / 0.358, unseen 18.79 / 0.472, PSNR@64 16.07, FVD 329 / 690, IDM 0.344; U-Net at 20k, seen 19.82 / 0.372, unseen 18.71 / 0.479, PSNR@64 15.90, FVD 265 / 531, IDM 0.326. The corrected DiT gains 0.9\,dB and 0.047 LPIPS on seen maps at 75k versus 90k.

\section{IDM details}\label{app:idm}
VPT-style bidirectional window of $K{=}8$ latents, shared convolutional encoder, 4-layer transformer, one 29-way head per transition, 11.4M parameters, trained on verified transitions from the training split. Held-out real windows: 82.9\% top-1, 54.8\% macro recall, 84.9\% on the movement collapse, majority 23.2\%. $K{=}2$ control: 70.1 / 40.6 / 72.7\%. Rollout scoring averages overlapping-window logits over the 64-frame sequence; the real reference uses the identical procedure on the same trajectories (0.847; majority 0.232). Rollout agreement: DiT s0 0.492, s1 0.476, U-Net 0.509, episode-bootstrap intervals overlapping. \todo{per-action-group table.}

\section{Blur control}\label{app:blur}
Gaussian blur at $\sigma \in \{0, 0.5, 1, 2, 4\}$\,px applied to each model's predictions before scoring at horizons 8 and 64 never raises PSNR: U-Net at 64, 16.02 at $\sigma{=}0$, 16.04 at 0.5, \todo{$\sigma{=}1$}, \todo{$\sigma{=}2$}, 15.78 at 4; DiT s0 at 64, 18.30 down to 17.78 at $\sigma{=}4$. The copy-last and VAE-ceiling reference values are 19.41\,dB seen / 18.50 unseen and 28.61 / 0.051 seen / 26.36 / 0.070 unseen. \todo{full grid for every row and both horizons, LPIPS included.}

\section{Bootstrap}\label{app:boot}
Episode-level bootstrap, 95\% intervals: resample episodes with replacement and recompute means over all windows or rollouts of the drawn episodes. Seen PSNR: DiT s0 21.06 [20.74, 21.38], U-Net 21.36 [21.01, 21.73]. Seen LPIPS: 0.311 [0.301, 0.320], 0.270 [0.261, 0.280]. Unseen PSNR: DiT s0 19.52 [19.32, 19.73], DiT s1 19.59 [19.40, 19.80], U-Net 19.14 [18.90, 19.39]. Horizon-64 PSNR, U-Net minus DiT s0: $-2.28$ [$-2.88$, $-1.70$]. \todo{number of resamples; intervals for every table cell.}
\end{document}
